\documentclass[11pt]{article}

\usepackage[a4paper,margin=1in]{geometry}
\usepackage[T1]{fontenc}
\usepackage[utf8]{inputenc}
\usepackage{lmodern}
\usepackage{amsmath, amssymb}
\usepackage{verbatim}
\usepackage{url}
\usepackage{graphicx}    % \resizebox for over-wide tables
\usepackage{rotating}    % sidewaystable: dedicated landscape page

% Section numbering depth
\setcounter{secnumdepth}{3}

\title{
  \textbf{enclawed}: A Configurable, Sector-Neutral Hardening Framework for
  Single-User AI Assistant Gateways
}

\author{%
  Alfredo Metere\\
  Metere Consulting, LLC.\\
  \texttt{alfredo.metere@metereconsulting.com}
}

\date{April 17, 2026}

\begin{document}

\maketitle

\begin{abstract}
We present \emph{enclawed}, a hard-fork hardening framework built on top of
the OpenClaw single-user personal artificial intelligence (AI) assistant
gateway. enclawed targets deployments that need attestable peer trust,
deny-by-default external connectivity, signed-module loading, and a
tamper-evident audit trail --- typically regulated industries such as
financial services, healthcare, defense contracting, regulated R\&D, and
government enclaves. The framework ships in two flavors: an \texttt{open}
flavor that preserves OpenClaw compatibility while still emitting audit,
classification, and data-loss-prevention (DLP) signals, and an
\texttt{enclaved} flavor that activates strict allowlists, Federal
Information Processing Standards (FIPS) cryptographic-module assertion,
mandatory module-manifest signature verification, and high-assurance
peer attestation for the Model Context Protocol (MCP). The classification ladder is fully data-driven: a deploying
organization selects from five built-in presets (generic, US-government,
healthcare, financial services, three-tier) or supplies its own JSON. We
accompany the implementation with a security review, a 207-case test
suite (149 unit tests, 58 adversarial pen-tests for tamper detection,
signature forgery, egress bypass, trust-root mutation, DLP evasion,
prompt injection, and code injection), real-time human-in-the-loop
control (per-agent pause / resume / stop and approval queues), a
memory-bounded secure transaction buffer with rollback (default cap
50\% of system RAM, configurable), a strict-mode TypeScript typecheck
of all 22 framework files, and a GitHub Actions workflow ready for
continuous integration. \textbf{enclawed is a hardening framework,
not an accredited compliance certification.} The deploying organization
remains responsible for hardware, validated cryptographic modules,
certified facilities, and assessor sign-off.
\end{abstract}

% =====================================================================
\section{Introduction}

Generative AI has crossed quickly from prototype into the operational
fabric of regulated industries. Financial-services firms use large
language models (LLMs) for research summarization touching material
non-public information (MNPI); healthcare systems apply them to
protected health information (PHI); defense contractors handle
controlled unclassified information (CUI) and International Traffic in
Arms Regulations (ITAR)-controlled materials; legal teams rely on
them for privileged-counsel work; pharmaceutical R\&D operates on
embargoed clinical-trial data. Each of these settings is governed by
frameworks that pre-date generative AI but apply to it directly --- the
Health Insurance Portability and Accountability Act (HIPAA) [8], the
European Union General Data Protection Regulation (GDPR) [7], the
Payment Card Industry Data Security Standard (PCI~DSS) [9], the
International Organization for Standardization / International
Electrotechnical Commission (ISO/IEC) 27001 [5], the American Institute
of Certified Public Accountants Service Organization Control 2 (SOC~2)
trust-services criteria [6], the National Institute of Standards and
Technology (NIST) Special Publication 800-53 [1], and NIST 800-171 [4]
--- and increasingly by AI-specific regimes such as the European Union
Artificial Intelligence Act (EU AI Act) [28], the NIST AI Risk
Management Framework (NIST AI RMF) [26], and ISO/IEC 42001 [27]
governing artificial-intelligence management systems.

Two observations motivate this work. First, the dominant open-source
multi-channel AI gateways were designed for the consumer
``personal-assistant'' market and ship with cloud-first defaults that
are fundamentally incompatible with regulated-enclave deployment.
Second, the security failure modes of generative AI are no longer
hypothetical: documented incidents include training-data extraction at
scale [31, 32], indirect prompt injection through retrieved content
[20, 21, 22], the 2023 Samsung corpus leak through interactive cloud
application programming interfaces (APIs), and the 2024 \emph{Moffatt
v.\ Air Canada} decision [35] establishing operator liability for AI
assistant outputs. The Open Web Application Security Project (OWASP)
Top~10 for LLM Applications [29] and MITRE Adversarial Threat Landscape
for Artificial-Intelligence Systems (ATLAS) [30] now codify the
relevant attack surface in industry-recognized form.

\subsection{The regulated-enterprise AI gap}

Personal AI gateways such as OpenClaw [25] exemplify the problem. The
upstream project's positioning is explicit: a single-user assistant
that ``answers you on the channels you already use,'' bundling 24+
chat platforms (WhatsApp, Telegram, Slack, Discord, Signal, iMessage,
Matrix, Microsoft Teams, and more) and 40+ cloud LLM providers (OpenAI,
Anthropic, Google, Mistral, Groq, AWS Bedrock, Azure OpenAI, OpenRouter,
Fireworks, etc.) by default. Its plugin model encourages
community-published modules installed via \texttt{npm} with no enforced
signature verification. Its policy posture is permissive: external
channels and providers are enabled the moment credentials are
configured. Audit, classification, and DLP are not part of the
framework's contract.

For a community user this is the right design. For a financial-services
trading floor handling MNPI, a hospital system processing PHI, a
defense contractor with CUI access, or any organization subject to ISO
27001 / SOC~2 / HIPAA / PCI / Federal Risk and Authorization Management
Program (FedRAMP) / U.S.\ Department of Defense (DoD) / U.S.\ Department
of Energy (DoE) controls, every one of those defaults becomes an
explicit policy violation:

\begin{itemize}
  \item \emph{External egress to consumer messaging platforms} violates
    boundary protection (NIST 800-53 SC-7 [1]) and information-flow
    control (AC-4).
  \item \emph{Cloud LLM API usage} routes regulated content through
    third parties whose data-processing terms typically permit
    retention and, until contractually disclaimed, training reuse ---
    in direct conflict with HIPAA business-associate constraints [8,
    §164.308(b)], GDPR Article~28 processor obligations [7], and
    most data-residency policies.
  \item \emph{Unsigned community plugins} expose the gateway to
    supply-chain compromise --- the threat the NIST Secure Software
    Development Framework (SSDF) [33] addresses through SI-7 / SR-3 /
    SR-4, and which the U.S.\ Cybersecurity and Infrastructure Security
    Agency (CISA) / National Security Agency (NSA) / Federal Bureau of
    Investigation (FBI) joint guidance on AI deployment [34] explicitly
    calls out as a top operational risk.
  \item \emph{Permissive defaults with no audit chain} make incident
    reconstruction impossible, conflicting with NIST 800-53 AU-2/3/9/10
    and HIPAA §164.312(b) audit-control requirements.
  \item \emph{No classification labels} means an output may carry
    PHI / MNPI / CUI to a recipient cleared for none of it, with the
    host silent at every step.
\end{itemize}

\subsection{Why a configuration-only fix is insufficient}

A natural response is ``just disable the cloud channels and providers,
add a logging plugin, and require approved modules.'' We argue this is
not enough, for three structural reasons.

\textbf{Defaults are policy.} A platform whose out-of-the-box state
contains 78 forbidden modules invites accidental enablement at every
upgrade, every onboarding, every plugin install. Removing the modules
from the source tree --- not just from the configuration --- is the
only durable defense, because it makes the unsafe state
\emph{unreachable} rather than merely \emph{unselected}. Compliance
posture must be a property of the binary, not a property of a
configuration file an operator can revert.

\textbf{The plugin-loading path is the trust boundary.} OpenClaw's
plugin loader does not verify signatures, does not consult any trust
root, and does not reason about clearance. Adding ``a logging plugin''
is itself a plugin --- it inherits the same unsigned, unverified
loading path. The trust boundary must be \emph{below} the plugin
layer, enforced by the host before any plugin code runs, and locked
against post-boot mutation. Retrofitting this into a project whose
contract assumes ambient plugin trust is a re-architecture, not a
configuration.

\textbf{Upstream cannot adopt these constraints.} OpenClaw's product
identity --- a multi-channel personal assistant --- is incompatible
with deny-by-default external connectivity, FIPS-required boot,
mandatory signature verification, and attested peer trust. Even if
upstream maintainers were willing to accept patches that imposed those
defaults, doing so would break their published contract with thousands
of community deployments. Forking is therefore not merely convenient
but structurally necessary: the regulated-enclave product must diverge
from the consumer-assistant product at the binary level. We say
explicitly that upstream OpenClaw \emph{will not} adopt these
constraints, not because the upstream maintainers are inattentive to
security --- they are not --- but because the constraints are
incompatible with the product they ship.

\subsection{Three design commitments}

This paper describes \emph{enclawed}, a hard fork of OpenClaw rebuilt
around three design commitments that address the gap above:

\begin{enumerate}
  \item \textbf{Always-on policy.} The classification framework
    activates unconditionally at process start, before any plugin or
    transit module is imported, so downstream code observes a guarded
    environment from the first instruction. There is no
    ``compatibility mode'' that disables the framework in production.
  \item \textbf{Two flavors.} An \texttt{open} flavor preserves
    OpenClaw behavior with the framework running in warn-only mode for
    development and non-regulated deployment; an \texttt{enclaved}
    flavor enforces strict deny-by-default policy, FIPS assertion,
    mandatory module-manifest signing, and attested peer trust.
    Selection is by a single environment variable read once at boot
    and then immutable for the process lifetime.
  \item \textbf{Data-driven classification.} The level ladder is not
    hardcoded; the deploying organization selects from five built-in
    presets (sector-neutral generic, U.S.-government, healthcare-HIPAA,
    financial-services, three-tier) or supplies its own JavaScript
    Object Notation (JSON) document. The framework is sector-neutral by
    default --- a U.S.-government preset and DoE Q-clearance and
    L-clearance templates exist as opt-ins, not as the canonical
    defaults.
\end{enumerate}

The framework is sector-neutral by default. Built-in presets cover
financial services (MNPI / privileged-counsel), healthcare (PHI /
sensitive PHI), a generic three-tier scheme, and full US-government
markings, with the canonical defaults industry-generic (\texttt{PUBLIC}
$<$ \texttt{INTERNAL} $<$ \texttt{CONFIDENTIAL} $<$ \texttt{RESTRICTED}
$<$ \texttt{RESTRICTED-PLUS}).

\subsection{Comparison with the broader landscape}
\label{sec:landscape}

OpenClaw is one of several open-source projects that mediate between a
human user and one or more language models. Adjacent offerings
include LibreChat [43], Open WebUI [44], AnythingLLM [45], Jan.ai
[46], Lobe Chat [47] --- multi-provider chat front-ends --- plus
agent-orchestration frameworks such as AutoGen [48] and the LangChain
ecosystem [49]. A separate cluster of AI-security tools targets
narrower slices of the problem: NeMo Guardrails [50] and Llama Guard
[51] focus on content moderation at the model boundary; Lakera Guard
[52] addresses prompt injection as a hosted scanning service.

We surveyed each of these against the security properties enclawed is
designed to deliver, summarized in Table~\ref{tab:landscape}. The
honest picture is that no other publicly-available offering, as
documented at the time of writing, combines mandatory module signing,
a classification lattice with a configurable scheme, tamper-evident
audit, real-time human-in-the-loop control, a memory-bounded rollback
buffer, and attested peer trust into a single cohesive framework that
runs entirely on locally-hosted inference. The closest competitors
each address one or two of these properties:

\begin{itemize}
  \item Local-LLM front-ends (Open WebUI, AnythingLLM, Jan.ai) deliver
    point (1) --- they avoid cloud egress by binding to local runtimes
    --- but ship without policy, classification, audit, or
    signed-module enforcement. Their plugin or knowledge-base ingest
    paths are unverified, which is the same supply-chain gap noted in
    §1.1 for OpenClaw.
  \item Multi-provider gateways (LibreChat, Lobe Chat) prioritize
    reach over restriction; cloud providers are first-class citizens
    and removal requires user effort that is not enforced. Their
    plugin / preset distribution is community-driven without
    cryptographic attestation.
  \item Agent frameworks (AutoGen, LangChain ecosystem) are libraries
    rather than gateways, and the security posture is whatever the
    embedding application provides. They have no notion of
    deployment-flavor strictness, no built-in egress allowlist, and
    no rollback on partial-failure of an agent's actions.
  \item Content-guardrail systems (NeMo Guardrails, Llama Guard)
    address output filtering --- a complement to, not a replacement
    for, classification-driven flow control. They do not enforce
    information-flow rules, do not produce a tamper-evident audit
    trail, and do not gate module loading.
  \item SaaS prompt-injection scanners (Lakera Guard) by design route
    content through an external service, which is incompatible with
    the no-egress requirement of regulated enclaves.
\end{itemize}

Where a competitor genuinely addresses one of enclawed's properties,
enclawed adopts the same primitive rather than reinventing it: the DLP
scanner and prompt-shield (§4.5, §8) are deliberately a thin
keyword-and-regex layer because content-classification depth is the
job of dedicated tools like Llama Guard, which a deploying
organization can wire in addition. enclawed's contribution is the
host-level chassis that any of these complementary tools can sit
inside, with the chassis itself taking responsibility for the
properties the tools cannot deliver alone (signed module loading,
classification lattice, deny-by-default policy, attested peer trust,
HITL controls, rollback).

\begin{table}[ht]
  \centering
  \scriptsize
  \setlength{\tabcolsep}{4pt}
  \renewcommand{\arraystretch}{1.2}
  \resizebox{\textwidth}{!}{%
  \begin{tabular}{|l|c|c|c|c|c|c|c|c|c|}
    \hline
    \textbf{Project} & \textbf{Local-only} & \textbf{Mand.\ signed} & \textbf{Class.} & \textbf{Tamper-} & \textbf{DLP} & \textbf{HITL} & \textbf{Tx} & \textbf{Attested} & \textbf{Sector-neutral} \\
    & \textbf{by default} & \textbf{modules} & \textbf{lattice} & \textbf{evident audit} & \textbf{at output} & \textbf{stop} & \textbf{rollback} & \textbf{peer trust} & \textbf{config scheme} \\
    \hline
    \textbf{enclawed (this work)} & yes (enclaved) & yes & yes & yes & yes & yes & yes & yes (MCP) & yes \\
    \hline
    OpenClaw [25]              & no & no & no & no & no & no & no & no & no \\
    \hline
    LibreChat [43]              & no & no & no & no & no & no & no & no & no \\
    \hline
    Open WebUI [44]             & yes & no & no & no & no & no & no & no & no \\
    \hline
    AnythingLLM [45]            & opt-in & no & no & no & no & no & no & no & no \\
    \hline
    Jan.ai [46]                 & yes & no & no & no & no & no & no & no & no \\
    \hline
    Lobe Chat [47]              & no & no & no & no & no & no & no & no & no \\
    \hline
    AutoGen [48]                & N/A (library) & no & no & no & no & partial & no & no & no \\
    \hline
    LangChain [49]              & N/A (library) & no & no & no & partial & partial & no & no & no \\
    \hline
    NeMo Guardrails [50]        & N/A (component) & no & no & no & yes & no & no & no & no \\
    \hline
    Llama Guard [51]            & N/A (model) & no & no & no & yes & no & no & no & no \\
    \hline
    Lakera Guard [52]           & no (SaaS) & no & no & no & yes & no & no & no & no \\
    \hline
  \end{tabular}%
  }%
  \caption{Landscape comparison. Properties documented as of the
    projects' publicly-stated features at time of writing; ``partial''
    means the property is achievable but not enforced by the project
    itself, and ``N/A'' means the project is a different kind of
    artifact (library, component, model) rather than a hostable
    gateway and the column does not directly apply. enclawed is the
    only entry that combines all nine properties as a host-level
    framework around which the others can be composed.}
  \label{tab:landscape}
\end{table}

\paragraph{Non-claim.} This work does not constitute a certification.
Real audits --- ISO 27001, SOC~2, the Health Information Trust Alliance
(HITRUST) common security framework, FedRAMP, DoD/DoE high-side
Authority to Operate (ATO), PCI DSS Report on Compliance (RoC) --- all
require accredited hardware, validated cryptographic modules (FIPS
140-3 [2] where the regime demands), certified facilities, evaluation
by a qualified assessor, and management sign-off. enclawed is a
code-level scaffold; the gaps a deploying organization must close are
catalogued in Section~\ref{sec:gaps}.

\paragraph{Contributions.}
\begin{itemize}
  \item A two-flavor classification framework with sector-neutral
    defaults and built-in presets for finance, healthcare, government,
    R\&D, and a minimal three-tier scheme.
  \item A zero-trust blockchained key broker (§4.7) for hybrid
    deployments that cannot achieve full enclave isolation: K-of-N
    quorum across independent providers, signed attestations, an
    Ed25519 broker-signed hash-chained ledger, and CONSENSUS plus
    THRESHOLD-XOR combining modes.
  \item A user-configurable classification scheme model that accepts
    arbitrary level vocabularies, validated against rank-contiguity and
    name-uniqueness invariants.
  \item An Ed25519 module-signing system with a per-signer
    clearance-approval trust root and boot-time pre-verification, plus
    an analogous attestation model for Model Context Protocol peers.
  \item A hash-chained, append-only audit log with concurrent-append
    serialization, deep payload sanitization against log injection, and
    independent third-party verification.
  \item An adversarial test suite of 58 cases targeting tamper
    detection, signature forgery, egress bypass, trust-root mutation,
    DLP evasion, prompt injection, and code injection, alongside 149
    unit tests, all passing on Node 22.
  \item A GitHub Actions workflow that runs the suites on push, PR, and
    weekly cron, plus a strict-mode TypeScript typecheck of all 22
    framework files.
\end{itemize}

% =====================================================================
\section{Background and Threat Model}

\subsection{Choice of access-control model}
\label{sec:bla-rationale}

Information security has produced a small library of formal access
control models, each emphasizing a different property. We surveyed the
candidates against three criteria: (i) the primary risk in our
deployments is unauthorized \emph{disclosure} of regulated data, not
modification or fraud; (ii) the policy must be \emph{mandatory} ---
imposed by the host, not configurable by users or modules; (iii) the
ladder of sensitivity levels must be \emph{reusable} across sectors
that already think in tier vocabularies (PHI, MNPI, CUI, etc.).
Table~\ref{tab:model-choice} summarizes the alternatives we
considered.

\begin{table}[ht]
  \centering
  \footnotesize
  \setlength{\tabcolsep}{4pt}
  \renewcommand{\arraystretch}{1.15}
  \begin{tabular}{|l|p{2.6cm}|p{4.5cm}|p{3.3cm}|}
    \hline
    \textbf{Model} & \textbf{Primary property} & \textbf{Why not the primary fit} & \textbf{Use in enclawed} \\
    \hline
    Bell-LaPadula [10] & confidentiality & --- & primary axis \\
    \hline
    Biba [36] & integrity & does not constrain disclosure & complementary, future work \\
    \hline
    Clark-Wilson [37] & integrity via well-formed transactions & overengineered for chat-mediated AI; assumes app-defined IVPs & inspires the rollback buffer (§4.6) \\
    \hline
    Brewer-Nash [38] & dynamic conflict-of-interest separation & sector-specific (financial); orthogonal to tier flow & easily layered via compartments \\
    \hline
    Harrison-Ruzzo-Ullman (HRU) access matrix & general expressiveness & undecidable safety; no enforced lattice & rejected \\
    \hline
    Role-based access control (RBAC) [39] & role-based authorization & no notion of data sensitivity; group$\to$resource only & rejected as primary axis \\
    \hline
    Attribute-based access control (ABAC) [40] & attribute expressions & expressive but harder to formally verify; risk of policy drift & adopted within compartments / releasability slots \\
    \hline
  \end{tabular}
  \caption{Access-control model survey. Bell-LaPadula is selected as
    the primary axis because the deployments enclawed targets are
    fundamentally confidentiality-driven (PHI, MNPI, CUI, embargoed
    research). Biba's integrity dual and Clark-Wilson's transaction
    discipline are complementary and inform §4.6 (rollback buffer);
    RBAC and ABAC are layered as compartment / releasability
    assignments rather than as the primary flow rule.}
  \label{tab:model-choice}
\end{table}

Concretely we chose Bell-LaPadula for five reasons:

\begin{enumerate}
  \item \emph{Confidentiality is the first-order risk.} Every regime in
    Section~\ref{sec:gaps}'s gap list is a confidentiality regime first
    (PHI under HIPAA [8], MNPI under U.S.\ Securities and Exchange
    Commission (SEC) rules, CUI under NIST 800-171 [4], classified
    information under NIST 800-53 [1]). Integrity and availability
    matter, but only after confidentiality is established.
  \item \emph{The lattice is mechanizable.} Dominance, least
    upper-bound, and reading/writing predicates are simple, total,
    decidable functions on integer ranks plus subset tests on
    compartment sets. They are easier to formally verify than the
    rule-evaluation of Clark-Wilson or the policy expressions of ABAC.
  \item \emph{It maps directly to existing vocabularies.} US-government
    classification, healthcare PHI tiers, and financial MNPI all
    naturally express as a rank ladder plus optional compartments. The
    same code consumes any of them via the configurable scheme of §6.
  \item \emph{It is mandatory by construction.} The lattice is a
    property of the host, evaluated on every label comparison. Modules
    cannot override it.
  \item \emph{Biba and Clark-Wilson compose orthogonally.} A future
    integrity layer can use Biba's no-read-down / no-write-up on a
    parallel axis without modifying the existing lattice; the
    transaction buffer (§4.6) already implements the well-formed
    transaction discipline of Clark-Wilson at the action layer.
\end{enumerate}

\subsection{Bell-LaPadula formal model}

Each subject and object carries a label $\langle \ell, C, R \rangle$
where $\ell$ is a numeric rank, $C$ a set of compartments, and $R$ a
set of releasability caveats. Label $a$ dominates label $b$ ($a
\succeq b$) iff $\ell_a \geq \ell_b$ and $C_b \subseteq C_a$. Reading
is allowed when the subject dominates the object (no-read-up); writing
is allowed when the object dominates the subject (no-write-down).
Combination yields the least upper bound. This lattice arithmetic is
implemented in \texttt{src/enclawed/classification.ts} and is
independent of the level vocabulary.

\subsection{Threat model}

We assume a single-tenant, single-user gateway running inside a
high-trust enclave. The user holds the deploying organization's highest
applicable trust tier. The system must:

\begin{enumerate}
  \item never egress to the public Internet or any external channel/provider;
  \item use only locally-hosted inference (e.g.\ Ollama [18], vLLM [16],
    LM Studio [19], SGLang [17], local NVIDIA Inference Microservice (NIM));
  \item refuse to render or echo data above the user's authorized tier;
  \item refuse to write data below its origin classification (no-write-down);
  \item produce a tamper-evident audit trail of every model interaction;
  \item encrypt every persistent artifact at rest with a key bound to
    organization-controlled hardware;
  \item block known sensitive-data markings, secrets, and personally
    identifiable information (PII) from leaving the gateway via any
    output channel;
  \item make any deviation a hard, loud failure --- not silent
    degradation.
\end{enumerate}

We map these requirements to multiple control frameworks simultaneously
(Table~\ref{tab:controls}): NIST 800-53 [1], ISO/IEC 27001 [5], NIST
Cybersecurity Framework (CSF) 2.0 [3], SOC~2 trust-services criteria
(TSC) [6], GDPR [7], the HIPAA Security Rule [8], PCI DSS [9], and the
Cybersecurity Maturity Model Certification (CMMC) layered on NIST
800-171 [4]. The deployment-side complement to this mapping --- the
controls the deploying organization must still deliver --- is
enumerated in Section~\ref{sec:gaps}.

\begin{sidewaystable}
  \centering
  \normalsize
  \setlength{\tabcolsep}{6pt}
  \renewcommand{\arraystretch}{1.3}
  \resizebox{\textheight}{!}{%
  \begin{tabular}{|l|l|l|l|l|l|l|l|}
    \hline
    \textbf{Control area} & \textbf{enclawed surface} & \textbf{NIST 800-53} & \textbf{ISO 27001} & \textbf{NIST CSF} & \textbf{SOC 2} & \textbf{GDPR} & \textbf{HIPAA} \\
    \hline
    Access enforcement & \texttt{classification.ts}, \texttt{policy.ts} & AC-3 & A.9.4 & PR.AC-4 & CC5.1, CC6.1 & Art.~32(1)(b) & 164.312(a)(1) \\
    \hline
    Information flow & \texttt{egress-guard.ts} & AC-4, SC-7 & A.13.1 & PR.PT-4 & CC6.6 & Art.~32(1)(b) & 164.312(e)(1) \\
    \hline
    Security attributes & \texttt{classification.ts} + scheme & AC-16 & A.8.2 & PR.DS-5 & CC1.4 & Art.~5(1)(f) & 164.308(a)(2) \\
    \hline
    Audit logging & \texttt{audit-log.ts} + \texttt{subsystem} patch & AU-2/3/9/10 & A.12.4 & DE.CM-1, DE.AE-3 & CC7.2, CC7.3 & Art.~30, 32(1)(d) & 164.312(b) \\
    \hline
    Crypto module & \texttt{crypto-fips.assertFipsMode()} & IA-7, SC-13 & A.10.1 & PR.DS-1 & CC6.1 & Art.~32(1)(a) & 164.312(a)(2)(iv) \\
    \hline
    Encryption at rest / in transit & \texttt{crypto-fips.ts} & SC-8/13/28 & A.10.1, A.13.2 & PR.DS-1/2 & CC6.1, CC6.7 & Art.~32(1)(a) & 164.312(e)(2)(ii) \\
    \hline
    Continuous monitoring & audit + subsystem + egress.onDeny & SI-4 & A.12.4 & DE.CM, DE.AE & CC7.2 & Art.~32(1)(d) & 164.308(a)(1)(ii)(D) \\
    \hline
    Data leakage / DLP & \texttt{dlp-scanner.ts} & SI-12 & A.13.2.3 & PR.DS-5 & CC6.7 & Art.~5(1)(f) & 164.312(e) \\
    \hline
    Memory hygiene & \texttt{zeroize.ts} & MP-6, SI-16 & A.8.3 & PR.IP-6 & CC6.5 & Art.~32(1)(a) & 164.310(d)(2)(i) \\
    \hline
    Least functionality & module deletions (Sec.~3.2) & CM-7 & A.9.1 & PR.AC-3 & CC6.3 & Art.~25 & 164.312(a)(1) \\
    \hline
    Module integrity & \texttt{module-signing} + \texttt{trust-root} & SI-7, CM-5, SR-3/4 & A.12.5, A.14.2.5 & PR.IP-1, PR.DS-6 & CC8.1 & Art.~32(1)(b) & 164.312(c)(1) \\
    \hline
    Peer attestation & \texttt{mcp-attested} & IA-3, IA-9, AC-3(11) & A.13.1.1 & PR.AC-7 & CC6.1 & Art.~32(1)(b) & 164.312(d) \\
    \hline
  \end{tabular}%
  }%
  \caption{Control mapping. Each row links a framework surface to the
  corresponding requirement in seven major frameworks; CMMC L2/L3
  inherits NIST 800-171 [4] and PCI DSS [9] adds Req.~3, 4, 7, 10, 11
  alongside NIST 800-53. The deploying organization picks the column
  that fits its regime. (Rendered on a dedicated landscape page; cells
  use natural-width \texttt{l} columns and the whole table is uniformly
  scaled to fill the page.)}
  \label{tab:controls}
\end{sidewaystable}

\subsection{Out of scope}

The framework does \emph{not} provide accredited cryptography
(FIPS 140-3 or NSA Type-1), an accredited cross-domain solution (CDS),
identity binding (Common Access Card (CAC), Personal Identity
Verification (PIV), Security Assertion Markup Language (SAML), or
OpenID Connect (OIDC)), kernel-level mandatory access control
(SELinux Multi-Level Security (MLS)), Write-Once-Read-Many (WORM)
audit shipping, or facility security. These remain the deploying
organization's responsibility. Section~\ref{sec:gaps} enumerates each.

% =====================================================================
\section{Architecture}

\subsection{Two flavors}

enclawed selects a flavor at boot via the \texttt{ENCLAWED\_FLAVOR}
environment variable. Table~\ref{tab:flavors} contrasts the two.

\begin{table}[t]
  \centering
  \small
  \begin{tabular}{|p{3.5cm}|p{5cm}|p{5cm}|}
    \hline
    \textbf{Aspect} & \textbf{open (default)} & \textbf{enclaved} \\
    \hline
    Channel/provider/host allowlists & not enforced & strict deny-by-default \\
    \hline
    Module signatures & warn-only & required, must verify \\
    \hline
    Signer-clearance approval & informational & required, exact match \\
    \hline
    FIPS assertion at boot & off (configurable) & on (default) \\
    \hline
    Trust root mutation post-boot & permitted & locked \\
    \hline
    \texttt{globalThis.fetch} reassignment & permitted & frozen non-configurable \\
    \hline
    Module without manifest & loaded with warning & rejected \\
    \hline
    MCP peer attestation & verified, fail $\to$ warn & verified, fail $\to$ deny \\
    \hline
  \end{tabular}
  \caption{Flavor matrix. Both flavors run the same framework code; the
  distinction is enforcement strictness.}
  \label{tab:flavors}
\end{table}

\subsection{Module set after deletion}

The fork removes 78 cloud module directories from the upstream
\texttt{extensions/} tree: 25 cloud channels (WhatsApp, Telegram, Slack,
Discord, Signal, iMessage, Matrix, MS Teams, etc.), 45 cloud LLM
providers (OpenAI, Anthropic, Google, Mistral, Groq, AWS Bedrock, Azure
OpenAI, OpenRouter, Fireworks, etc.), and 8 external search/browser/
webhook modules (Brave, DuckDuckGo, Exa, Firecrawl, SearXNG, Tavily,
Browser, Webhooks). 35 modules remain, all local or local-capable
(Ollama, vLLM, LM Studio, SGLang, local NVIDIA NIM, ComfyUI, the
media/memory/speech cores, openshell, qa-channel, phone-control,
device-pair, diagnostics-otel for local-only collectors). One new
module ships: \texttt{mcp-attested}, described in Section~\ref{sec:mcp}.

\subsection{Runtime singleton and bootstrap}

Activation is a single call at the top of \texttt{src/entry.ts}, before
any plugin or transit module is imported:

\begin{verbatim}
const { bootstrapEnclawed } = await import("./enclawed/bootstrap.js");
await bootstrapEnclawed();
\end{verbatim}

\texttt{bootstrapEnclawed()} reads the flavor and classification scheme
from the environment, loads the policy, opens the audit log, installs
the egress guard (with \texttt{freeze:true} in \texttt{enclaved}),
pre-verifies every module manifest under \texttt{extensions/}, locks the
trust root in \texttt{enclaved}, and stores the runtime on a
\texttt{Symbol.for("enclawed.runtime")} slot of \texttt{globalThis} so
downstream patches can consult it without creating import cycles. The
boot record (with active scheme id, flavor, FIPS state, and allowlist
contents) is the genesis hash-chain entry.

% =====================================================================
\section{Implementation}

The framework is delivered in two parallel surfaces: a TypeScript (TS)
twin under \texttt{src/enclawed/} that bundles with the upstream
OpenClaw build, and a canonical reference under \texttt{enclawed/src/}
written in plain Node ECMAScript Modules (\texttt{.mjs}) with zero
runtime dependencies. The
canonical surface runs under \texttt{node --test} without
\texttt{pnpm install}, which keeps the standalone security suite cheap to
exercise and trivial to audit. Table~\ref{tab:loc} summarizes file counts
and lines of code per tree.

\paragraph{Repository layout: two git repositories with a submodule
boundary.} The project is split into two \emph{independent git
repositories}, each with its own license and remote, at the repository
boundary:

\begin{itemize}
  \item \texttt{enclawed-oss} --- open-source, MIT-licensed,
    self-contained, public-bound git repository. Houses the OpenClaw
    fork base, the framework primitives (classification, policy,
    audit, DLP, crypto wrapper, zeroize, module signing, trust root,
    HITL, transaction buffer, prompt shield, egress guard), the
    canonical \texttt{.mjs} reference, the standalone test suite,
    this paper, and all open / local-capable \texttt{extensions/}
    (Ollama, vLLM, LM Studio, SGLang, ComfyUI, qa-channel, openshell,
    the media / memory / speech cores, etc.).
  \item \texttt{enclawed-enclaved} --- closed-source extensions,
    proprietary all-rights-reserved, private-bound git repository.
    Houses the zero-trust blockchained key broker (§4.7), the bundled
    \texttt{mcp-attested} reference module (§\ref{sec:mcp}), and the
    reference classified-profile config.
    \textbf{Consumes \texttt{enclawed-oss} as a git submodule}
    (\texttt{.gitmodules} pins a specific OSS commit; for sibling
    GitHub hosting the URL is \texttt{../enclawed-oss}, which resolves
    to the public OSS repository at the same organization).
    \texttt{git clone --recurse-submodules} pulls both halves in one
    step; closed-tree imports use relative paths into the populated
    submodule (e.g.\ \texttt{../../enclawed-oss/src/enclawed/audit-log.js}
    from \texttt{src/enclawed-secure/}).
\end{itemize}

The dependency direction is strict:
\texttt{enclawed-enclaved} $\to$ \texttt{enclawed-oss}, never the
reverse. Removing the entire closed repository must leave the OSS
repository fully buildable and testable. Each repository ships its
own GitHub Actions workflow: the OSS workflow runs the OSS suite
alone; the closed-tree workflow uses \texttt{actions/checkout@v4}
with \texttt{submodules: recursive} to fetch the OSS half in one
step, then runs the closed-tree suite plus a manifest-signature
verification step over every signed module the closed tree bundles.

This is a deliberate \emph{open core} arrangement: the parts that
benefit most from broad community review --- the classification model,
the audit-log primitive, the module-signing toolkit, the prompt
shield --- are open. The parts that carry distribution-controlled
value to the deploying organization --- the zero-trust key broker that
mediates between the host and external custodians, the attested-MCP
reference module with its real signed manifest, the curated
classified-profile config --- ship under separate licensing in a
private repository that pins the OSS submodule by commit hash.

\begin{table}[t]
  \centering
  \small
  \begin{tabular}{|l|r|r|}
    \hline
    \textbf{Surface (tree)} & \textbf{Files} & \textbf{LOC} \\
    \hline
    \multicolumn{3}{|l|}{\emph{enclawed-oss/ (open source, MIT)}} \\
    \hline
    TypeScript framework (\texttt{src/enclawed/*.ts}, integration shims) & 22 & $\sim$2{,}500 \\
    \hline
    Canonical \texttt{.mjs} reference (\texttt{enclawed/src/}) & 17 & $\sim$1{,}450 \\
    \hline
    Standalone unit + pen tests (\texttt{enclawed/test/}) & 19 & $\sim$1{,}700 \\
    \hline
    Upstream patches (4 files) & 4 & $+$82 \\
    \quad \texttt{src/entry.ts} & & $+$10 \\
    \quad \texttt{src/plugins/channel-validation.ts} & & $+$27 \\
    \quad \texttt{src/plugins/provider-validation.ts} & & $+$27 \\
    \quad \texttt{src/logging/subsystem.ts} & & $+$18 \\
    \hline
    \multicolumn{3}{|l|}{\emph{enclawed-enclaved/ (closed source, depends on OSS)}} \\
    \hline
    Zero-trust key broker (TS + MJS + tests) & 3 & $\sim$880 \\
    \hline
    \texttt{mcp-attested} module (TS + tests + signed manifest) & 5 & $\sim$305 \\
    \hline
    Reference classified-profile config & 1 & 41 \\
    \hline
  \end{tabular}
  \caption{Code footprint of the two repositories. The OSS repository
  is self-contained; the closed repository consumes the OSS repository
  as a git submodule and reaches into it via relative paths.}
  \label{tab:loc}
\end{table}

\subsection{Hash-chained audit log}

Every record is canonicalized as JSON with sorted keys, prepended with
the previous record's SHA-256, and hashed. Three properties matter:

\begin{itemize}
  \item \emph{Tamper detection.} Editing a record in place breaks the
    chain at the next record's \texttt{prevHash} comparison, which
    \texttt{verifyChain()} reports as \texttt{brokenAt}.
  \item \emph{Concurrent-write safety.} \texttt{append()} chains every
    write onto an internal \texttt{Promise} queue so two simultaneous
    callers cannot both consume the same \texttt{lastHash}. Without this
    guard, parallel emitters (e.g.\ the egress-guard \texttt{onDeny}
    callback racing with \texttt{createSubsystemLogger}) would produce a
    self-inconsistent chain.
  \item \emph{Log-injection neutralization.} A \texttt{deepSanitize}
    pass strips C0 control characters from every string in
    \texttt{type}, \texttt{actor}, \texttt{level}, and the entire
    payload tree, replacing them with U+FFFD. This both prevents an
    attacker from spoofing a fake JSONL record by embedding a newline in
    a payload string, and ensures the on-disk bytes match the bytes
    that were hashed. \texttt{\_\_proto\_\_}, \texttt{constructor}, and
    \texttt{prototype} keys are filtered out before canonicalization so
    a payload object cannot smuggle prototype-pollution shapes into the
    audit hash.
\end{itemize}

\subsection{Egress guard}

\texttt{installEgressGuard()} replaces \texttt{globalThis.fetch} with a
wrapper that resolves the target hostname via the Web Hypertext
Application Technology Working Group (WHATWG) URL parser (yielding
case-normalized hostnames, ignoring the \texttt{Host} header, and
rejecting embedded credentials) and checks against an allowlist. In
the \texttt{enclaved} flavor the property is set via
\texttt{Object.defineProperty} with \texttt{writable:false} and
\texttt{configurable:false}, making subsequent reassignment throw under
ECMAScript Modules (ESM) strict mode. The guard explicitly does not
protect against raw \texttt{net.Socket}, \texttt{dgram}, Domain Name
System (DNS) resolutions inside native modules, or child processes;
the deploying organization must layer kernel-level egress controls
(\texttt{nftables}, extended Berkeley Packet Filter (eBPF), network
namespaces, optional one-way diode) on top.

\subsection{Configurable classification scheme}

A scheme is a JSON document conforming to:

\begin{verbatim}
{
  "id": "acme-2026",
  "description": "ACME Corp internal data classification policy v3.2",
  "levels": [
    { "rank": 0, "canonicalName": "Public",        "aliases": ["P"] },
    { "rank": 1, "canonicalName": "Internal",      "aliases": ["I"] },
    { "rank": 2, "canonicalName": "Customer Data", "aliases": [] },
    { "rank": 3, "canonicalName": "Privileged",    "aliases": ["legal"] }
  ],
  "validCompartments":  ["FINANCE", "ENG", "LEGAL"],
  "validReleasability": ["NDA", "EYES_ONLY"]
}
\end{verbatim}

\texttt{parseClassificationScheme()} enforces three invariants: ranks
must be contiguous starting at zero; every name (canonical and alias)
must be unique across the whole scheme after case-insensitive
normalization; non-empty levels list. The active scheme governs
\texttt{format()}, \texttt{parse()}, \texttt{makeLabel()} range
validation, manifest \texttt{clearance} validation, and required-tier
checks at the MCP attestation layer. Five presets ship out of the box:
\texttt{enclawed-default} (six tiers; generic + US-gov merged),
\texttt{us-government} (UNCLASSIFIED to TOP SECRET//SCI),
\texttt{healthcare-hipaa} (PHI / Sensitive-PHI / Research-Embargoed),
\texttt{financial-services} (MNPI / Privileged-Counsel),
\texttt{generic-3-tier} (Public / Internal / Restricted).

\subsection{DLP scanner}

A regex-based scanner with three pattern families:
\emph{sensitive-data markings} (industry banners, U.S.-government
classification banners, DoE Restricted Data, Sensitive Compartmented
Information (SCI) codewords, distribution caveats); \emph{cloud and
vendor secrets} (Amazon Web Services (AWS), Google Cloud Platform (GCP)
service account, Microsoft Azure storage key, GitHub, GitLab, OpenAI,
Anthropic, Slack, Stripe, JSON Web Tokens (JWTs), Privacy-Enhanced
Mail (PEM)-encoded private keys); and \emph{international PII} (email
address, ITU-T E.164 phone number, credit-card primary account number
(PAN), International Bank Account Number (IBAN), U.S.\ Social Security
Number (SSN)). Input length is capped at 1\,MiB by default to bound
CPU exposure to regular-expression denial-of-service (ReDoS)-shaped
patterns; callers may pass \texttt{onOversize:'truncate'} to scan the
prefix. Severity levels (\texttt{low}, \texttt{medium}, \texttt{high},
\texttt{critical}) drive a \texttt{redact()} entry point that replaces
matches above a threshold with \texttt{[REDACTED]}. The scanner is a
backstop, not a primary control: regex DLP cannot detect paraphrased
sensitive content, optical-character-recognition-only content, or any
pattern not in its catalog.

\subsection{Human-in-the-loop control}
\label{sec:hitl}

\texttt{src/enclawed/hitl.ts} provides per-agent sessions with a
state machine \texttt{PENDING $\to$ RUNNING $\rightleftarrows$ PAUSED
$\to$ \{STOPPED, COMPLETED, FAILED\}}, a real-time event stream on the
controller (\texttt{EventEmitter}-based, with both a typed channel per
event name and a generic \texttt{event} channel), and an approval
queue. Each \texttt{AgentSession} exposes:

\begin{itemize}
  \item \texttt{pause()}, \texttt{resume()}, \texttt{stop(reason)},
    \texttt{complete()}, \texttt{fail(reason)} --- explicit operator
    controls. \texttt{stop} is final; the state machine forbids
    transitions out of any terminal state.
  \item \texttt{checkpoint()} --- a cooperative cancellation point
    the agent code awaits between actions. It throws
    \texttt{AgentStoppedError} if the session has been stopped, and
    blocks until \texttt{resume()} if the session is paused. We chose
    cooperative over preemptive cancellation deliberately: Saltzer and
    Schroeder's \emph{complete mediation} principle [41] argues for
    explicit checkpoints over forced thread termination, which in a
    single-threaded async runtime can leave shared state in an
    inconsistent half-mutated form.
  \item \texttt{proposeAction(actionType, payload)} --- if the
    action type belongs to the session's approval-required set, the
    call awaits a human decision (allow / deny) routed through the
    controller's approval queue. Deny raises
    \texttt{ApprovalDeniedError}; if the agent is stopped while the
    call is awaiting the human, it instead raises
    \texttt{AgentStoppedError} (the approval is no longer relevant).
\end{itemize}

The \texttt{HitlController} maintains the session registry, the
approval queue, and a \texttt{stopAll(reason)} mass-halt operation.
Every state transition and every approval request/resolution is
emitted as an event and (when an \texttt{AuditLogger} is wired)
appended to the hash-chained audit log under the
\texttt{hitl.}$\langle$\texttt{type}$\rangle$ namespace. A separate
operator UI (CLI, web, or remote) subscribes to the event stream to
render live status and to call \texttt{resolveApproval(id, decision)}
when a human approves or denies a queued action. Sub-second latency
is the design goal; the controller does no I/O on the hot path.

\subsection{Secure transaction buffer with rollback}
\label{sec:tx}

The framework treats every reversible action an agent takes as a
\emph{transaction}, in the sense of Gray's classical formulation
[42]. The deploying organization registers each action's inverse at
the moment of execution; the buffer (\texttt{src/enclawed/transaction-buffer.ts})
holds the resulting transaction record so it can be rolled back later
if a human, an automated guardrail, or a failed downstream operation
demands undo.

\paragraph{Memory-bounded rollback window.} The buffer is sized as a
configurable percentage of system RAM (default 50\%, configurable via
the \texttt{ramPercent} constructor option, with an absolute
\texttt{maxBytes} override). The size cap matches the eviction policy
of Clark-Wilson [37]'s well-formed transaction discipline applied to a
finite resource: as long as the action sits in the buffer it can be
inverted; once evicted, it has been implicitly committed. A new record
that would push the buffer past the cap auto-commits the oldest
records FIFO until enough room is freed.

\paragraph{Tamper-evident chain.} Each transaction is hash-chained to
its predecessor with SHA-256, in the same shape as the audit log of
§4.1. \texttt{verifyChain()} performs an independent walk so an
operator (or a separate auditor process) can detect silent tampering
of the in-memory record set. Every record, rollback, commit, and
eviction is emitted to the audit log when one is wired.

\paragraph{Rollback semantics.} \texttt{rollback(n)} pops the n most
recent records and runs each \texttt{inverse} in LIFO order. An
inverse that throws is reported in the result's \texttt{errors} array
and an audit \texttt{transaction.rollback-failed} record is appended,
but rollback continues so a single broken inverse cannot block
restoration of the rest of the window. After rollback the chain is
re-anchored to the new tail so subsequent records continue cleanly.

\subsection{Zero-trust blockchained key broker}
\label{sec:zerotrust}

The deployments enclawed targets prefer full enclave isolation, but
some are not free to choose: a healthcare system may be required to
escrow keys in a regulator-approved managed Key Management Service
(KMS); a defense contractor may operate at CMMC Level 2 [4] with a
FedRAMP cloud KMS; a financial-services firm may use multiple cloud
hardware security modules (HSMs) across regions for availability. In
these settings the key custodian sits outside the host's trust
boundary --- the broker must work without trusting any single provider.

\texttt{src/enclawed/zero-trust-key-broker.ts} implements a host-side
broker that:

\begin{itemize}
  \item Treats every \texttt{KeyProvider} as untrusted. Each response
    is a \texttt{providerId}-pinned, Ed25519-signed
    \texttt{Attestation} carrying \texttt{(keyId, ts, keyHash, payload,
    nonce)}; an unsigned, malformed, replayed (\texttt{keyId}
    mismatch), cross-claimed (\texttt{providerId} mismatch), or
    self-inconsistent (\texttt{keyHash $\neq$ SHA-256(payload)})
    attestation is excluded from the quorum count.
  \item Requires a configurable K-of-N quorum of valid agreeing
    attestations before it returns any key material to the host.
  \item Records every operation --- success and quorum-not-met
    failures --- in a hash-chained, broker-signed ledger
    (\texttt{KeyChainLedger}). The ledger is the ``blockchain'':
    single-writer, append-only, with each block containing the
    operation, the ordered list of providers queried, fingerprints of
    every attestation considered, the SHA-256 of the derived key
    bytes, the previous block's hash, and an Ed25519 broker signature.
    \texttt{verifyChain(brokerPublicKeyPem)} performs an independent
    walk so any external auditor with the broker's public key plus the
    providers' public keys can validate the entire history without
    trusting the broker either.
  \item Supports two operating modes:
    \begin{enumerate}
      \item \textbf{Consensus mode} --- each provider returns the same
        wrapped key blob (e.g.\ they all decrypt the same envelope
        from a shared root). The broker groups attestations by their
        declared \texttt{keyHash} and accepts the largest group that
        reaches quorum. A dishonest provider returning a different
        value is excluded silently from the winning group.
      \item \textbf{Threshold-XOR mode} --- each provider returns one
        XOR share of the key. The broker requires all N shares
        (degenerate K=N quorum) and reconstructs the master key by
        XOR. No single provider ever holds the full key; collusion of
        N$-$1 providers is insufficient.
    \end{enumerate}
\end{itemize}

The architecture deliberately mirrors the audit log of §4.1 and the
transaction buffer of §4.6: the same hash-chain primitive
\cite{HaberStornetta1991, Crosby2009} is reused at three layers
(audit, rollback, key operations), so an external auditor only needs
to learn one verification routine.

\paragraph{Limitations.} Threshold-XOR is N-of-N, not K-of-N. A future
Shamir secret-sharing [53] backend can supply true K-of-N
confidentiality of the master key while preserving the same
\texttt{KeyProvider} interface; the broker is structured to accept
additional combiners. Like the audit log, the ledger detects
in-place tamper but not tail truncation, so production deployments
must replicate it to a separate auditor in real time. Provider-side
key storage is the provider's problem; the broker only verifies what
each provider says it returned. Byzantine-fault analysis of the
quorum threshold follows the classical bounds of
Lamport-Shostak-Pease [54] (with $f$ Byzantine providers, the broker
needs $K \geq f+1$ valid attestations to ensure at least one honest
provider is in the winning group). Concretely: at $N=5$ providers,
$K=3$ tolerates one Byzantine provider; at $N=7$, $K=4$ tolerates
two; the deploying organization picks the K/N pair that matches its
threat model. NIST SP 800-152 [55] is a useful reference for the
operational constraints on managed key custodians.

% =====================================================================
\section{Module Signing and Trust Root}

Each module ships an \texttt{enclawed.module.json} alongside its
\texttt{package.json}:

\begin{verbatim}
{ "v": 1,
  "id": "mcp-attested",
  "publisher": "enclawed reference",
  "version": "0.1.0",
  "clearance": "restricted-plus",
  "capabilities": ["mcp-client", "tool"],
  "signerKeyId": "enclawed-attested-reference-2026",
  "signature": "<base64 Ed25519 signature>" }
\end{verbatim}

The signature is an Ed25519 [13] signature over a stable canonical
encoding of the manifest body excluding \texttt{signature} itself, so
signing is idempotent. The trust root binds each public key to a list
of clearance tiers it is approved to attest. Verification passes when
(i) the manifest's \texttt{signerKeyId} resolves to a signer in the
trust root, (ii) the signer is approved for the manifest's declared
clearance, (iii) the signature verifies against the canonical bytes,
and (iv) any caller-required clearance is met. In the \texttt{enclaved}
flavor, the bootstrap calls \texttt{lockTrustRoot()} after the
deploying organization's signers are loaded, after which any
\texttt{setTrustRoot()} or \texttt{resetTrustRoot()} call throws
\texttt{TrustRootLockedError}.

The module loader is two-staged:
\texttt{preloadModuleDecisions(rootDir)} walks the modules directory
once at boot, parses every manifest, runs \texttt{checkModule()}
against the active flavor and trust root, and stashes the per-id
decisions on the runtime singleton; the synchronous validation
chokepoints (channel- and provider-validation, both patched into
upstream OpenClaw) read the cached decisions without performing async
I/O on a hot path. Modules denied at preload are rejected when their
channels or providers subsequently try to register, with a
corresponding \texttt{module.deny.channel} or
\texttt{module.deny.provider} audit record.

% =====================================================================
\section{Configurable Classification --- Worked Examples}

Table~\ref{tab:schemes} shows the five built-in schemes side by side.

\begin{table}[t]
  \centering
  \footnotesize
  \setlength{\tabcolsep}{3pt}
  \resizebox{\textwidth}{!}{%
  \begin{tabular}{|c|l|l|l|l|l|}
    \hline
    \textbf{Rank} & \textbf{default} & \textbf{us-government} & \textbf{healthcare-hipaa} & \textbf{financial-services} & \textbf{3-tier} \\
    \hline
    0 & PUBLIC          & UNCLASSIFIED    & PUBLIC             & PUBLIC             & PUBLIC \\
    \hline
    1 & INTERNAL        & CUI             & INTERNAL           & INTERNAL           & INTERNAL \\
    \hline
    2 & CONFIDENTIAL    & CONFIDENTIAL    & PHI                & CONFIDENTIAL       & RESTRICTED \\
    \hline
    3 & RESTRICTED      & SECRET          & SENSITIVE-PHI      & MNPI               & --- \\
    \hline
    4 & RESTRICTED-PLUS & TOP SECRET      & RESEARCH-EMB.\     & PRIVILEGED-COUNSEL & --- \\
    \hline
    5 & SCI             & TOP SECRET//SCI & ---                & ---                & --- \\
    \hline
  \end{tabular}%
  }%
  \caption{Built-in classification schemes (column headers are scheme
    ids; rank-4 of \texttt{healthcare-hipaa} is RESEARCH-EMBARGOED,
    abbreviated for table fit). Aliases such as \texttt{q-cleared}
    $\equiv$ rank 4 of \texttt{us-government} are accepted by
    \texttt{parse()} and by manifest \texttt{clearance} validation.}
  \label{tab:schemes}
\end{table}

A deploying organization selects a scheme via:

\begin{verbatim}
ENCLAWED_FLAVOR=enclaved \
ENCLAWED_CLASSIFICATION_SCHEME=healthcare-hipaa \
ENCLAWED_AUDIT_PATH=/var/log/enclawed/audit.jsonl \
node openclaw.mjs gateway run --bind loopback
\end{verbatim}

For a custom scheme, the same env var is set to a JSON file path;
\texttt{loadSchemeByName} accepts an optional \texttt{allowedDirs}
parameter that rejects paths outside a vetted configuration directory.

% =====================================================================
\section{Attested MCP Module}
\label{sec:mcp}

\texttt{mcp-attested} is a reference Model Context Protocol [15] client
that gates every connection on a remote peer's signed clearance
attestation. The remote server publishes a JSON document at
\texttt{/.well-known/} with the same shape as
\texttt{enclawed.module.json}; the client fetches it, runs it through
\texttt{checkModule()} with the caller's required clearance, and
proceeds only on success. The capability \texttt{mcp-server} must be
declared. A typical refusal sequence:

\begin{verbatim}
const c = new QClearedMcpClient({
  requiredClearance: "restricted-plus",
});
const r = await c.connect("https://mcp-internal.acme.lan:8443/");
// In enclaved flavor:
//   { ok: false, reason: "signer not approved..." }
// In open flavor:
//   { ok: false, reason: "...(open flavor: warn-only)" }
\end{verbatim}

The client is sector-neutral: \texttt{requiredClearance} can be any
name in the active scheme, so a healthcare deployment can require
\texttt{sensitive-phi} and a U.S.-government deployment can require
\texttt{q-cleared} using the same code path.

% =====================================================================
\section{Prompt Injection Defenses}

\texttt{src/enclawed/prompt-shield.ts} provides defensive helpers for
text that flows from untrusted sources (channel input, retrieved
document, tool output) into a model prompt. The framework does not
promise complete prompt-injection defense --- that is an open research
problem [20, 21, 22] --- but it neutralizes the most common
silent-confusion vectors so they surface as visible content rather than
control signals to the model.

\texttt{sanitizeForPrompt(text)} performs five operations in order:

\begin{enumerate}
  \item \texttt{stripControlChars} replaces ASCII C0 control characters
    (the byte range 0x00--0x1F excluding TAB) with the Unicode
    replacement character U+FFFD.
  \item \texttt{stripBidi} removes the Unicode bidirectional-override
    range U+202A--U+202E and U+2066--U+2069.
  \item \texttt{stripZeroWidth} removes the zero-width characters
    U+200B--U+200D, U+2060, and U+FEFF.
  \item \texttt{neutralizeRoleBoundaries} prefixes any line starting
    with \texttt{system:} / \texttt{assistant:} / \texttt{user:} /
    \texttt{tool:} / \texttt{function:} (with optional leading
    whitespace or Markdown emphasis) with a literal
    \texttt{[USER-CONTENT]} marker so the model sees them as quoted
    rather than as a fresh role boundary.
  \item \texttt{neutralizeFences} inserts a zero-width character before
    any line consisting of triple backticks so the closing fence cannot
    break out of a quoted code block the host wraps untrusted input in.
\end{enumerate}

\texttt{detectInjection(text)} returns finding identifiers without
modifying the text, including a heuristic for ALL-CAPS imperative
override phrases (``IGNORE ALL PREVIOUS INSTRUCTIONS'', ``OVERRIDE the
above messages'', etc.).

% =====================================================================
\section{Evaluation}

\subsection{Test methodology}

The framework ships with two test surfaces. The canonical \texttt{.mjs}
suite uses Node's built-in \texttt{node:test} runner [24] and depends
only on \texttt{node:crypto}, allowing verification without any
\texttt{npm install}. A vitest twin under
\texttt{src/enclawed/integration.test.ts} exercises the same surface in
the upstream OpenClaw \texttt{pnpm test} pipeline.

\subsection{Unit test coverage}

Table~\ref{tab:unit} enumerates the per-module unit-test counts. All
149 unit tests pass on Node 22.

\begin{table}[t]
  \centering
  \small
  \begin{tabular}{|l|r|p{6cm}|}
    \hline
    \textbf{Module} & \textbf{Tests} & \textbf{Coverage focus} \\
    \hline
    \texttt{classification.test.mjs}        & 13 & lattice arithmetic, format/parse round-trip, immutability \\
    \hline
    \texttt{classification-scheme.test.mjs} & 11 & all 5 presets round-trip, custom JSON, validation invariants \\
    \hline
    \texttt{dlp-scanner.test.mjs}           & 11 & all 3 pattern families, redact correctness \\
    \hline
    \texttt{policy.test.mjs}                & 9  & enclaved/open defaults, allowlist enforcement, frozen \\
    \hline
    \texttt{module-loader.test.mjs}         & 8  & 8 deny paths and 1 allow path under both flavors \\
    \hline
    \texttt{crypto-fips.test.mjs}           & 10 & AES-GCM round-trip, AAD round-trip (string + Buffer), AAD mismatch rejection, FIPS gate \\
    \hline
    \texttt{zeroize.test.mjs}               & 7  & buffer fill, withSecret on success and throw \\
    \hline
    \texttt{egress-guard.test.mjs}          & 6  & block, allow, lifecycle \\
    \hline
    \texttt{module-manifest.test.mjs}       & 6  & valid parse, invariant rejection, hash stability \\
    \hline
    \texttt{flavor.test.mjs}                & 5  & alias parsing, env override, default \\
    \hline
    \texttt{audit-log.test.mjs}             & 4  & append + verify, tamper, reopen, hash determinism \\
    \hline
    \texttt{module-signing.test.mjs}        & 4  & Ed25519 round-trip, wrong-key, tampered, malformed \\
    \hline
    \texttt{hitl.test.mjs}                  & 14 & state machine, cooperative checkpoint, pause/resume/stop, approval queue, audit integration \\
    \hline
    \texttt{transaction-buffer.test.mjs}    & 18 & ramPercent sizing, hash chain, LIFO rollback, partial-failure rollback, eviction, audit integration \\
    \hline
    \texttt{zero-trust-key-broker.test.mjs} & 19 & attestation sign/verify, quorum success and failure, signature forgery exclusion, providerId / keyId / payload-hash mismatch exclusion, threshold-XOR reconstruction, ledger verifyChain (clean / tampered / wrong key), audit integration \\
    \hline
    \texttt{mcp-attested/test/server-clearance.test.mjs} & 4 & generic and US-gov vocab, deny-on-low, open warn-only \\
    \hline
    \textbf{Total unit + module} & \textbf{149} & \\
    \hline
  \end{tabular}
  \caption{Unit test inventory.}
  \label{tab:unit}
\end{table}

\subsection{Adversarial pen-test coverage}

Table~\ref{tab:pen} enumerates the per-file pen-test counts. All 58
pen-tests pass on Node 22.

\begin{table}[t]
  \centering
  \small
  \begin{tabular}{|l|r|p{7cm}|}
    \hline
    \textbf{File} & \textbf{Tests} & \textbf{Adversarial focus} \\
    \hline
    \texttt{audit-log.pentest.mjs}        & 6  & in-place edit, tail truncation (documented limit), record reorder, newline log-injection, \texttt{\_\_proto\_\_} pollution, 100 concurrent appends \\
    \hline
    \texttt{signature-forgery.pentest.mjs} & 7  & wrong-key, signerKeyId swap, signature replay across modules, downgrade attack, capability injection, malformed signature, open-flavor warning capture \\
    \hline
    \texttt{egress-bypass.pentest.mjs}    & 8  & hostname case normalization, Host-header spoofing, embedded credentials, IPv4-in-IPv6, file:/data: schemes, fetch reassign with/without freeze (latter via subprocess), malformed URLs \\
    \hline
    \texttt{trust-root-and-scheme.pentest.mjs} & 10 & setTrustRoot after lock, \texttt{\_\_proto\_\_}/constructor pollution, duplicate normalized names, non-contiguous ranks, negative rank, empty levels, non-string canonicalName, missing id, type checks \\
    \hline
    \texttt{dlp-evasion.pentest.mjs}      & 8  & oversize cap, ReDoS bound, zero-width camouflage, whitespace variants, redact correctness, PEM camouflage gap, null/undefined safety \\
    \hline
    \texttt{prompt-injection.pentest.mjs} & 11 & 5 spoofed roles, bidi overrides, zero-width camouflage, control chars, code-fence breakout, imperative-override phrases, multi-vector composite, idempotency \\
    \hline
    \texttt{code-injection.pentest.mjs}   & 8  & shell-metachar manifest fields, JS-eval clearance rejected, function-shaped JSON treated as data, channel id path-traversal denied, \texttt{\_\_proto\_\_} provider id, allowedDirs allowlist, malformed JSON error context, deep-JSON safety \\
    \hline
    \textbf{Total pen-tests} & \textbf{58} & \\
    \hline
  \end{tabular}
  \caption{Adversarial test inventory. Each row covers a distinct
  attack family; cases that document a known limitation are noted as
  ``documented'' or ``gap'' so the deploying organization knows where
  defense-in-depth is required.}
  \label{tab:pen}
\end{table}

\subsection{Bugs found and fixed during construction}

The unit and pen suites caught the following real defects during
development; each was fixed prior to publication:

\begin{enumerate}
  \item \texttt{crypto-fips.deriveKey} ---
    \texttt{ERR\_CRYPTO\_INVALID\_SCRYPT\_PARAMS} at boundary. Lifted
    \texttt{maxmem} explicitly to 64\,MiB.
  \item \texttt{classification.makeLabel} --- compartments backed by
    \texttt{Set} were mutable through \texttt{.add()} despite
    \texttt{Object.freeze}. Switched to deduplicated, sorted, frozen
    arrays.
  \item \texttt{classification.parse} --- auto-promoting
    \texttt{TOP\_SECRET} to \texttt{TOP\_SECRET\_SCI} on any
    non-releasability segment broke \texttt{format}/\texttt{parse}
    round-tripping. Removed; SCI is a separately representable head.
  \item \texttt{audit-log.append} --- concurrent calls raced on
    \texttt{lastHash}, producing a broken chain. Fixed with internal
    Promise queue.
  \item \texttt{audit-log} log-injection --- payload strings containing
    raw newlines could spoof a fake JSONL record. Fixed with
    \texttt{deepSanitize} pass before record construction.
  \item \texttt{trust-root} --- post-boot \texttt{setTrustRoot} could
    substitute attacker-controlled signers. Fixed with
    \texttt{lockTrustRoot()} called by the bootstrap in
    \texttt{enclaved}.
  \item \texttt{dlp-scanner} --- unbounded input could trigger
    pathological regex CPU. Added 1\,MiB cap with throw or truncate.
  \item \texttt{egress-guard} --- \texttt{globalThis.fetch} reassignment
    could bypass the guard. Added \texttt{freeze:true}, which sets the
    property via \texttt{Object.defineProperty} with
    \texttt{writable:false} and \texttt{configurable:false}.
  \item Scheme parser --- accepted a non-string \texttt{canonicalName}
    via \texttt{String(...)} coercion. Now requires a real string with
    TypeError.
  \item Scheme loader --- unwrapped \texttt{JSON.parse} surfaced a
    bare \texttt{SyntaxError}. Wrapped with file-path context; added
    \texttt{allowedDirs} path-traversal allowlist.
\end{enumerate}

\subsection{Aggregate result}

All 207 tests (149 unit + 58 pen) pass on Node 22 in approximately
0.4\,seconds. All 22 framework TypeScript files type-check clean under
\texttt{tsc --strict --noEmit}. The continuous-integration (CI)
workflow at
\texttt{.github/workflows/enclawed-security.yml} runs both suites on
push, pull request (PR), and weekly cron, plus a strict-mode typecheck
and a boot-time signature verification of every shipped module
manifest.

% =====================================================================
\section{Related Work}

\paragraph{Multi-channel AI assistants.} OpenClaw [25] is the upstream
personal-AI gateway from which enclawed is forked. LiteLLM [19] provides
a multi-provider proxy but does not address classification or
attestation. Local-inference engines such as Ollama [18], vLLM [16],
LM Studio [19], and SGLang [17] are the inference targets that survive
the deletion catalog of Section 3.2.

\paragraph{Information-flow control.} The Bell-LaPadula model [10]
originated the no-read-up / no-write-down formulation enclawed
implements. Operating-system mandatory-access-control (MAC)
implementations such as SELinux [11] are the kernel-layer counterpart
that enclawed assumes the deploying organization will configure
separately.

\paragraph{Prompt-injection research.} Greshake et al.\ [20] formalized
indirect prompt injection; Perez and Ribeiro [21] catalogued early
attack patterns; Schulhoff et al.\ [22] demonstrated competition-grade
jailbreaks. enclawed's \texttt{prompt-shield} module addresses the
data-sanitization layer of this problem only; the broader research
direction remains open.

\paragraph{Attested peer trust.} Transport Layer Security (TLS) client
certificates and the Secure Production Identity Framework for Everyone
(SPIFFE) and its SPIRE runtime [23] solve identity attestation in
service meshes. enclawed's \texttt{mcp-attested} module borrows the
same shape (signed manifest plus trust root) but applies it
specifically to the Model Context Protocol [15] where peer attestation
is not yet standardized.

\paragraph{Audit logs.} Tamper-evident logging via hash chains is a
classical technique [12, 14]. enclawed implements a single-machine
variant; production deployments should layer WORM storage or off-host
shipping (e.g.\ to a separate audit aggregator) on top.

% =====================================================================
\section{Limitations and Gaps the Deploying Organization Owns}
\label{sec:gaps}

The framework cannot, by itself, deliver an accredited deployment. The
following gaps are explicitly out of scope and must be filled by the
deploying organization's security and compliance program:

\begin{itemize}
  \item \emph{Identity binding.} \texttt{canRead} enforces
    Bell-LaPadula, but no identity layer ships. Integrate the
    organization's identity provider (IdP) --- using SAML, OIDC, or a
    smart card --- and bind a clearance label to every session.
  \item \emph{OS-level mandatory access control.} Defense in depth
    requires SELinux MLS or equivalent enforcing the same lattice
    outside the JavaScript process.
  \item \emph{Cross-trust-zone transfer.} The framework forbids egress;
    a real deployment needs an accredited cross-zone control (a data
    diode, a file-transfer appliance, a manual review queue, or a
    vendor CDS).
  \item \emph{Validated cryptographic module.} The deployment must link
    against a regime-appropriate validated module (FIPS 140-3 for U.S.\
    federal / DoD / DoE / FedRAMP; Common Criteria (CC) Evaluation
    Assurance Level (EAL) evaluations for some EU and Asia-Pacific
    (APAC) regimes).
  \item \emph{Key management.} Replace passphrase parameters with
    HSM-backed key references (Public-Key Cryptography Standards \#11
    (PKCS\#11), the Key Management Interoperability Protocol (KMIP),
    or the vendor's software development kit (SDK)).
  \item \emph{At-rest coverage.} Wrap every persistent artifact ---
    credentials, sessions, model weights, swap, core dumps --- not just
    the credential planner.
  \item \emph{Audit log durability.} Ship records to WORM media or an
    isolated audit aggregator; the on-disk hash chain alone cannot
    detect trailing truncation.
  \item \emph{Memory hygiene.} JavaScript strings cannot be zeroized;
    audit every secret-handling path. Disable swap and core dumps.
  \item \emph{Plugin allowlist by signature, not id.} Bind each allowed
    plugin to a signed manifest hash and reject any drift.
  \item \emph{Cascade work after module deletion.} The 78 deleted module
    directories leave dangling references in upstream code (legacy
    migrations, doctor repair paths, hard-coded fallbacks). The
    deploying organization must run \texttt{pnpm tsgo \&\& pnpm test
    \&\& pnpm build} to enumerate and fix them.
\end{itemize}

% =====================================================================
\section{Conclusion}

enclawed packages a coherent set of hardening primitives ---
classification labels, deny-by-default policy, egress allowlist,
hash-chained audit, DLP scanning, FIPS-gated at-rest crypto, secret
zeroization, Ed25519 module signing with a clearance-aware trust root,
and prompt-injection sanitization --- behind a single one-call bootstrap
that runs before any plugin loads. The two-flavor design preserves
upstream compatibility for community deployments while supporting strict
enforcement in regulated settings. The classification ladder is
data-driven, accommodating any sector's vocabulary through five
built-in presets or a custom JSON scheme. The 207-test suite (149 unit +
58 adversarial) and the CI workflow let a deploying organization
continuously verify that the security primitives still hold against the
attack patterns enumerated in Section 9.

The framework is not a substitute for accredited cryptography, hardware,
facilities, or assessor sign-off; the gaps in Section~\ref{sec:gaps}
remain the deploying organization's responsibility. What enclawed
provides is an explicit, well-tested, and continuously-verifiable
scaffold that the organization's compliance program can build on rather
than reinvent.

% =====================================================================
\section*{Availability and Licensing}

The project is split into two \emph{independent git repositories} with
different licenses, hosting visibilities, and CI pipelines. The closed
repository consumes the open repository as a git submodule.

\textbf{Open core (\texttt{enclawed-oss}, public repository):} The
framework primitives --- 22 TypeScript files plus the parallel
\texttt{.mjs} reference, the unit and adversarial pen-test suites, the
upstream OpenClaw fork base, the open / local-capable extensions
(Ollama, vLLM, LM Studio, etc.), this paper, and the OSS-tree CI
workflow --- are released under the \textbf{MIT License} and hosted at
\url{https://github.com/metereconsulting/enclawed} [56]. Copyright (c)
2026 Metere Consulting, LLC.; copyright (c) 2025 Peter Steinberger
(upstream OpenClaw). MIT permits commercial use, modification,
redistribution, and incorporation into proprietary projects, subject
only to preservation of the copyright notice and license text. To
clone:
\begin{verbatim}
git clone https://github.com/metereconsulting/enclawed
\end{verbatim}

\textbf{Closed extensions (\texttt{enclawed-enclaved}, private
repository):} The zero-trust blockchained key broker (§4.7), the
bundled \texttt{mcp-attested} reference module (§\ref{sec:mcp}), and
the reference classified-profile config are released under a
\textbf{proprietary, all-rights-reserved license}. Copyright (c) 2026
Metere Consulting, LLC. Use, modification, redistribution, or
preparation of derivative works requires a separate written license
agreement; see \texttt{LICENSE} in the closed repository and the
contact information in its \texttt{NOTICE.md}.

The closed repository \texttt{enclawed-enclaved} consumes
\texttt{enclawed-oss} as a git submodule (recorded in
\texttt{.gitmodules}, pinned to a specific OSS commit). The submodule
is initialized with \texttt{git clone --recurse-submodules} or
\texttt{git submodule update --init --recursive}. After
initialisation the OSS tree appears at \texttt{enclawed-oss/} inside
the closed repository, and closed-tree imports resolve into it via
relative paths. The submodule retains its own MIT \texttt{LICENSE} and
\texttt{NOTICE.md} inside the closed repository tree; the proprietary
license of the closed repository does not extend to anything inside
the submodule, and the MIT terms of the submodule do not extend to
anything outside it (MIT does not impose copyleft on dependents).
Removing the entire closed repository leaves the OSS repository fully
buildable, testable, and publishable.

A top-level \texttt{LICENSE.md} at the development workspace root
(alongside both repository checkouts) summarizes the
two-license posture, lists the conditions under which the arrangement
is legally well-formed (copyright ownership of the closed tree, no
copyleft contamination, contributor agreements for the OSS tree,
trademark search, patent and export-control review), and explicitly
disclaims that the licensing files were prepared without attorney
review. The publisher (Metere Consulting, LLC.) is solely responsible
for obtaining qualified IP-counsel verification before relying on this
posture in production.

% =====================================================================
\begin{thebibliography}{99}

\bibitem{NIST80053r5}
National Institute of Standards and Technology.
\emph{Security and Privacy Controls for Information Systems and
Organizations}, NIST Special Publication 800-53 Rev. 5, 2020.

\bibitem{NISTFIPS1403}
National Institute of Standards and Technology.
\emph{Security Requirements for Cryptographic Modules}, FIPS 140-3,
2019.

\bibitem{NISTCSF2}
National Institute of Standards and Technology.
\emph{The NIST Cybersecurity Framework (CSF) 2.0}, NIST CSWP 29, 2024.

\bibitem{NIST800171}
National Institute of Standards and Technology.
\emph{Protecting Controlled Unclassified Information in Nonfederal
Systems and Organizations}, NIST SP 800-171 Rev. 2, 2020.

\bibitem{ISO27001}
International Organization for Standardization.
\emph{ISO/IEC 27001:2022 --- Information security, cybersecurity and
privacy protection --- Information security management systems ---
Requirements}, 2022.

\bibitem{SOC2}
American Institute of Certified Public Accountants.
\emph{Trust Services Criteria for Security, Availability, Processing
Integrity, Confidentiality, and Privacy}, 2022.

\bibitem{GDPR}
European Parliament and Council of the European Union.
\emph{Regulation (EU) 2016/679 (General Data Protection Regulation)},
2016.

\bibitem{HIPAASecurityRule}
U.S. Department of Health and Human Services.
\emph{HIPAA Security Rule, 45 CFR Part 160 and Subparts A and C of Part
164}.

\bibitem{PCIDSSv4}
PCI Security Standards Council.
\emph{Payment Card Industry Data Security Standard, v4.0}, 2022.

\bibitem{BellLaPadula1976}
D.~E.~Bell and L.~J.~LaPadula.
\emph{Secure Computer System: Unified Exposition and Multics
Interpretation}, MITRE Technical Report MTR-2997, 1976.

\bibitem{Loscocco2001}
P.~Loscocco and S.~Smalley.
Integrating Flexible Support for Security Policies into the Linux
Operating System.
In \emph{Proceedings of the FREENIX Track: 2001 USENIX Annual
Technical Conference}, 2001.

\bibitem{HaberStornetta1991}
S.~Haber and W.~S.~Stornetta.
How to Time-Stamp a Digital Document.
\emph{Journal of Cryptology}, 3(2), 1991.

\bibitem{RFC8032}
S.~Josefsson and I.~Liusvaara.
\emph{Edwards-Curve Digital Signature Algorithm (EdDSA)}, IETF RFC
8032, 2017.

\bibitem{Crosby2009}
S.~A.~Crosby and D.~S.~Wallach.
Efficient Data Structures for Tamper-Evident Logging.
In \emph{18th USENIX Security Symposium}, 2009.

\bibitem{AnthropicMCP2024}
Anthropic.
\emph{Model Context Protocol Specification}, 2024.
\url{https://modelcontextprotocol.io/}.

\bibitem{Kwon2023vLLM}
W.~Kwon, Z.~Li, S.~Zhuang, Y.~Sheng, L.~Zheng, C.~H.~Yu, J.~Gonzalez,
H.~Zhang, and I.~Stoica.
Efficient Memory Management for Large Language Model Serving with
PagedAttention.
In \emph{Proceedings of the 29th Symposium on Operating Systems
Principles (SOSP)}, 2023.

\bibitem{Zheng2023SGLang}
L.~Zheng et al.
\emph{SGLang: Efficient Execution of Structured Language Model Programs}.
arXiv:2312.07104, 2023.

\bibitem{Ollama}
Ollama. \url{https://github.com/ollama/ollama}.

\bibitem{LMStudio}
LM Studio. \url{https://lmstudio.ai/}.

\bibitem{Greshake2023}
K.~Greshake, S.~Abdelnabi, S.~Mishra, C.~Endres, T.~Holz, and M.~Fritz.
Not What You've Signed Up For: Compromising Real-World LLM-Integrated
Applications with Indirect Prompt Injection.
In \emph{16th ACM Workshop on Artificial Intelligence and Security
(AISec)}, 2023.

\bibitem{Perez2022}
F.~Perez and I.~Ribeiro.
Ignore Previous Prompt: Attack Techniques For Language Models.
arXiv:2211.09527, 2022.

\bibitem{Schulhoff2023}
S.~Schulhoff et al.
Ignore This Title and HackAPrompt: Exposing Systemic Vulnerabilities
of LLMs through a Global Scale Prompt Hacking Competition.
In \emph{Proceedings of EMNLP}, 2023.

\bibitem{SPIFFE}
SPIFFE: Secure Production Identity Framework for Everyone.
\url{https://spiffe.io/}.

\bibitem{NodeTest}
Node.js Foundation.
\emph{Node.js node:test module}, 2024.
\url{https://nodejs.org/api/test.html}.

\bibitem{OpenClawRepo}
OpenClaw upstream repository.
\url{https://github.com/openclaw/openclaw}.

\bibitem{NIST_AI_RMF}
National Institute of Standards and Technology.
\emph{Artificial Intelligence Risk Management Framework (AI RMF 1.0)},
NIST AI 100-1, January 2023.

\bibitem{ISO42001}
International Organization for Standardization.
\emph{ISO/IEC 42001:2023 --- Information technology --- Artificial
intelligence --- Management system}, 2023.

\bibitem{EUAIAct}
European Parliament and Council of the European Union.
\emph{Regulation (EU) 2024/1689 laying down harmonised rules on
artificial intelligence (Artificial Intelligence Act)}, June 2024.

\bibitem{OWASP_LLM}
OWASP Foundation.
\emph{OWASP Top 10 for Large Language Model Applications}, 2025
edition.
\url{https://owasp.org/www-project-top-10-for-large-language-model-applications/}.

\bibitem{MITREATLAS}
The MITRE Corporation.
\emph{Adversarial Threat Landscape for Artificial-Intelligence Systems
(ATLAS)}.
\url{https://atlas.mitre.org/}.

\bibitem{Carlini2021}
N.~Carlini, F.~Tramer, E.~Wallace, M.~Jagielski, A.~Herbert-Voss,
K.~Lee, A.~Roberts, T.~Brown, D.~Song, U.~Erlingsson, A.~Oprea, and
C.~Raffel.
Extracting Training Data from Large Language Models.
In \emph{30th USENIX Security Symposium}, 2021.

\bibitem{Nasr2023}
M.~Nasr, N.~Carlini, J.~Hayase, M.~Jagielski, A.~F.~Cooper, D.~Ippolito,
C.~A.~Choquette-Choo, E.~Wallace, F.~Tramer, and K.~Lee.
\emph{Scalable Extraction of Training Data from (Production) Language
Models}.
arXiv:2311.17035, 2023.

\bibitem{NISTSSDF}
National Institute of Standards and Technology.
\emph{Secure Software Development Framework (SSDF) Version 1.1},
NIST SP 800-218, February 2022.

\bibitem{CISA_AI}
Cybersecurity and Infrastructure Security Agency, U.S. National
Security Agency, and U.S. Federal Bureau of Investigation.
\emph{Engaging with Artificial Intelligence: Joint Guidance for Secure
Deployment}, 2024.

\bibitem{MoffattAirCanada}
Civil Resolution Tribunal of British Columbia.
\emph{Moffatt v.\ Air Canada}, 2024 BCCRT 149, February 2024.

\bibitem{Biba1977}
K.~J.~Biba.
\emph{Integrity Considerations for Secure Computer Systems}, MITRE
Technical Report MTR-3153, 1977.

\bibitem{ClarkWilson1987}
D.~D.~Clark and D.~R.~Wilson.
A Comparison of Commercial and Military Computer Security Policies.
In \emph{Proceedings of the IEEE Symposium on Security and Privacy},
1987.

\bibitem{BrewerNash1989}
D.~F.~C.~Brewer and M.~J.~Nash.
The Chinese Wall Security Policy.
In \emph{Proceedings of the IEEE Symposium on Security and Privacy},
1989.

\bibitem{Sandhu1996}
R.~S.~Sandhu, E.~J.~Coyne, H.~L.~Feinstein, and C.~E.~Youman.
Role-Based Access Control Models.
\emph{IEEE Computer}, 29(2):38--47, February 1996.

\bibitem{NIST800162}
V.~C.~Hu, D.~Ferraiolo, R.~Kuhn, A.~Schnitzer, K.~Sandlin, R.~Miller,
and K.~Scarfone.
\emph{Guide to Attribute Based Access Control (ABAC) Definition and
Considerations}, NIST Special Publication 800-162, January 2014.

\bibitem{SaltzerSchroeder1975}
J.~H.~Saltzer and M.~D.~Schroeder.
The Protection of Information in Computer Systems.
\emph{Proceedings of the IEEE}, 63(9):1278--1308, September 1975.

\bibitem{Gray1981}
J.~Gray.
The Transaction Concept: Virtues and Limitations.
In \emph{Proceedings of the 7th International Conference on Very Large
Data Bases (VLDB)}, 1981.

\bibitem{LibreChat}
D.~Avila and contributors.
LibreChat --- enhanced ChatGPT clone with multi-provider, plugins, and
agents. \url{https://github.com/danny-avila/LibreChat}.

\bibitem{OpenWebUI}
Open WebUI Project.
Open WebUI --- user-friendly self-hosted AI interface (formerly
ollama-webui). \url{https://github.com/open-webui/open-webui}.

\bibitem{AnythingLLM}
Mintplex Labs.
AnythingLLM --- the all-in-one AI app for any LLM with full RAG and
agent capabilities. \url{https://github.com/Mintplex-Labs/anything-llm}.

\bibitem{JanAI}
Jan.ai Project.
Jan --- bring AI to your desktop. Open-source ChatGPT alternative,
runs 100\% offline. \url{https://github.com/janhq/jan}.

\bibitem{LobeChat}
LobeHub.
Lobe Chat --- open-source modern-design AI chat framework.
\url{https://github.com/lobehub/lobe-chat}.

\bibitem{AutoGen}
Q.~Wu, G.~Bansal, J.~Zhang, Y.~Wu, B.~Li, E.~Zhu, L.~Jiang, X.~Zhang,
S.~Zhang, J.~Liu, A.~H.~Awadallah, R.~W.~White, D.~Burger, and
C.~Wang.
AutoGen: Enabling Next-Gen LLM Applications via Multi-Agent
Conversation.
arXiv:2308.08155, 2023.

\bibitem{LangChain}
H.~Chase and the LangChain community.
LangChain --- building applications with LLMs through composability.
\url{https://github.com/langchain-ai/langchain}.

\bibitem{NeMoGuardrails}
T.~Rebedea, R.~Dinu, M.~Sreedhar, C.~Parisien, and J.~Cohen.
NeMo Guardrails: A Toolkit for Controllable and Safe LLM Applications
with Programmable Rails.
arXiv:2310.10501, 2023.

\bibitem{LlamaGuard}
H.~Inan, K.~Upasani, J.~Chi, R.~Rungta, K.~Iyer, Y.~Mao, M.~Tontchev,
Q.~Hu, B.~Fuller, D.~Testuggine, and M.~Khabsa.
Llama Guard: LLM-based Input-Output Safeguard for Human-AI
Conversations.
arXiv:2312.06674, 2023.

\bibitem{LakeraGuard}
Lakera AI.
Lakera Guard --- protect your LLM applications from prompt injection,
PII leaks, and data poisoning.
\url{https://www.lakera.ai/lakera-guard}.

\bibitem{Shamir1979}
A.~Shamir.
How to Share a Secret.
\emph{Communications of the ACM}, 22(11):612--613, November 1979.

\bibitem{LamportShostakPease1982}
L.~Lamport, R.~Shostak, and M.~Pease.
The Byzantine Generals Problem.
\emph{ACM Transactions on Programming Languages and Systems},
4(3):382--401, July 1982.

\bibitem{NIST800152}
National Institute of Standards and Technology.
\emph{A Profile for U.S. Federal Cryptographic Key Management Systems},
NIST Special Publication 800-152, October 2015.

\bibitem{EnclawedOSSRepo}
A.~Metere (Metere Consulting, LLC.).
\emph{enclawed --- open-source half}. Public repository.
\url{https://github.com/metereconsulting/enclawed}.

\end{thebibliography}

\end{document}
